\documentclass[11pt]{article}
\usepackage{amsmath,amssymb,amsthm}
\usepackage[margin=1in]{geometry}

\newcommand{\R}{\mathbb{R}}
\newcommand{\vmip}{v^{\mathrm{MIP}}}
\newcommand{\ssum}{\textstyle\sum}
\newtheorem{proposition}{Proposition}
\theoremstyle{remark}
\newtheorem{remark}[proposition]{Remark}

\title{Measuring $\varepsilon(\chi)$: a mixed-integer fractional program\\ and its Dinkelbach solution}
\date{}

\begin{document}
\maketitle

\section{The quantity to be measured}

Fix an allocation $\chi\in\R^{n}$ of the grand coalition's value among the members
$N=\{1,\dots,n\}$. Its worst per-capita excess is
\begin{equation}
\varepsilon(\chi)\;=\;\max_{\emptyset\neq S\subsetneq N}\
\frac{\vmip(S)-\ssum_{j\in S}\chi_j}{|S|},
\label{eq:pcexcess}
\end{equation}
so that $\varepsilon(\chi)\le 0$ certifies $\chi\in\mathcal{C}(v)$ and, more generally,
$\varepsilon(\chi)\le\bar\varepsilon$ certifies $\chi\in\mathcal{C}_{\bar\varepsilon}(v)$.

Two things make \eqref{eq:pcexcess} hard. The feasible set has $2^{n}-2$ elements, and
each evaluation of the numerator requires $\vmip(S)$, the optimal value of the coalition's
unit-commitment dispatch. Enumeration is therefore out of reach, and the maximisation
must be carried out implicitly.

\section{A single mixed-integer fractional program}

Introduce $z\in\{0,1\}^{n}$ with $S(z)=\{j: z_j=1\}$, and let $x$ collect the dispatch
variables, continuous and binary, of the community model. Write $X(z)$ for the dispatch
feasible set in which every variable belonging to a member $j$ with $z_j=0$ is forced to
zero, and $\pi^{\top}x$ for the dispatch objective, so that
\begin{equation}
\vmip(S(z))\;=\;\max_{x\in X(z)}\ \pi^{\top}x .
\label{eq:inner}
\end{equation}

The nesting in \eqref{eq:pcexcess} collapses. Because the numerator contains
$\vmip(S)$ with a positive sign, and $\vmip(S)$ is itself a maximum, the outer maximisation
over $S$ and the inner maximisation over $x$ are maxima of the same expression and may be
merged into one. Putting
\begin{equation}
\mathcal{N}(z,x)\;=\;\pi^{\top}x-\ssum_{j\in N}\chi_j z_j,
\qquad
\mathcal{D}(z)\;=\;\ssum_{j\in N} z_j,
\label{eq:numden}
\end{equation}
and
\begin{equation}
\Omega\;=\;\Big\{(z,x)\ :\ z\in\{0,1\}^{n},\ \ 1\le\ssum_{j\in N}z_j\le n-1,\ \ x\in X(z)\Big\},
\label{eq:omega}
\end{equation}
we obtain
\begin{equation}
\varepsilon(\chi)\;=\;\max_{(z,x)\in\Omega}\ \frac{\mathcal{N}(z,x)}{\mathcal{D}(z)} .
\label{eq:mifp}
\end{equation}
The bounds in \eqref{eq:omega} exclude $S=\emptyset$ and $S=N$ and, incidentally, guarantee
$\mathcal{D}(z)\ge 1>0$ throughout, which is what the fractional programme requires.

Problem \eqref{eq:mifp} is a \emph{mixed-integer linear fractional program}. Both
$\mathcal{N}$ and $\mathcal{D}$ are affine, the deactivation of a non-member is expressed by
big-$M$ inequalities rather than by products $z_j x_j$, and the nonlinear conversion
efficiencies are represented by a disaggregated piecewise-linear formulation with segment
binaries. No product of two variables appears anywhere, so $\Omega$ is a mixed-integer
\emph{linear} set. This is not a matter of convenience: it is what makes the certificate of
Section~\ref{sec:conv} available, since a global maximum over $\Omega$ is what
branch-and-bound delivers and what a nonconvex mixed-integer nonlinear formulation would not.

\section{The Dinkelbach transform}

Following Dinkelbach, associate with \eqref{eq:mifp} the parametric problem
\begin{equation}
F(\lambda)\;=\;\max_{(z,x)\in\Omega}\ \Big\{\mathcal{N}(z,x)-\lambda\,\mathcal{D}(z)\Big\},
\qquad \lambda\in\R .
\label{eq:Flambda}
\end{equation}

\begin{proposition}\label{prop:F}
$F$ is convex, strictly decreasing, and has a unique root, which is $\varepsilon(\chi)$:
\[
F(\lambda)>0\iff\lambda<\varepsilon(\chi),\qquad
F(\lambda)=0\iff\lambda=\varepsilon(\chi),\qquad
F(\lambda)<0\iff\lambda>\varepsilon(\chi).
\]
\end{proposition}

\begin{proof}
For fixed $(z,x)$ the map $\lambda\mapsto\mathcal{N}(z,x)-\lambda\mathcal{D}(z)$ is affine
with slope $-\mathcal{D}(z)\le-1$. $F$ is the pointwise maximum of these affine functions
over the finite set $\Omega$, hence convex and strictly decreasing. For the root, note that
$F(\lambda)\ge 0$ iff some $(z,x)\in\Omega$ has
$\mathcal{N}(z,x)/\mathcal{D}(z)\ge\lambda$, i.e.\ iff $\varepsilon(\chi)\ge\lambda$, and
$F(\lambda)\le 0$ iff every $(z,x)$ has
$\mathcal{N}(z,x)/\mathcal{D}(z)\le\lambda$, i.e.\ iff $\varepsilon(\chi)\le\lambda$.
\end{proof}

Computing $\varepsilon(\chi)$ is therefore root-finding on $F$, and Newton's method on a
convex decreasing function is exactly Dinkelbach's iteration:

\begin{enumerate}
\item Choose any $(z_0,x_0)\in\Omega$ and set
      $\lambda_1=\mathcal{N}(z_0,x_0)/\mathcal{D}(z_0)$.
\item For $k=1,2,\dots$ solve \eqref{eq:Flambda} at $\lambda_k$ to global optimality,
      obtaining a maximiser $(z_k,x_k)$ and the value $F(\lambda_k)$.
\item If $F(\lambda_k)\le 0$, stop: $\lambda_k=\varepsilon(\chi)$.
      Otherwise set $\lambda_{k+1}=\mathcal{N}(z_k,x_k)/\mathcal{D}(z_k)$ and repeat.
\end{enumerate}

\section{The parametric subproblem is the original separation}

The step that makes the scheme practical is that \eqref{eq:Flambda} is not a new model.
Substituting \eqref{eq:numden},
\begin{equation}
\mathcal{N}(z,x)-\lambda\,\mathcal{D}(z)
\;=\;\pi^{\top}x-\ssum_{j\in N}\chi_j z_j-\lambda\ssum_{j\in N}z_j
\;=\;\pi^{\top}x-\ssum_{j\in N}\big(\chi_j+\lambda\big)z_j .
\label{eq:shift}
\end{equation}
Because $|S|$ is linear in $z$, the term $\lambda\mathcal{D}(z)$ is absorbed into the
objective coefficients of the very binaries that define $S$. The parametric problem at
$\lambda$ is thus the raw-excess separation
\[
\max_{S\subseteq N}\ \Big\{\vmip(S)-\ssum_{j\in S}\chi'_j\Big\}
\qquad\text{with}\qquad \chi'_j=\chi_j+\lambda ,
\]
i.e.\ the same mixed-integer program with the same variables, constraints and big-$M$
structure, differing only in $n$ objective coefficients, each shifted by a common $\lambda$.
Every Dinkelbach iteration therefore costs one separation solve and no more.

\begin{remark}
A denominator nonlinear in $z$ would break this. The reduction in \eqref{eq:shift} is
available precisely because the cardinality $|S|=\ssum_j z_j$ is affine in the incidence
vector.
\end{remark}

\section{Convergence and what an early stop yields}\label{sec:conv}

\begin{proposition}
If every parametric problem \eqref{eq:Flambda} is solved to global optimality, the sequence
$(\lambda_k)$ is strictly increasing until it stops, satisfies $\lambda_k\le\varepsilon(\chi)$
for all $k$, and reaches $\varepsilon(\chi)$ after finitely many iterations.
\end{proposition}

\begin{proof}
Each $\lambda_{k+1}$ is the ratio attained by the feasible point $(z_k,x_k)$, so
$\lambda_{k+1}\le\varepsilon(\chi)$ by \eqref{eq:mifp}. If the method does not stop at
iteration $k$ then $F(\lambda_k)>0$, that is
$\mathcal{N}(z_k,x_k)-\lambda_k\mathcal{D}(z_k)>0$, and dividing by
$\mathcal{D}(z_k)\ge 1$ gives $\lambda_{k+1}>\lambda_k$. Since $\Omega$ contains finitely
many distinct values of $z$, and hence the ratios in \eqref{eq:mifp} take finitely many
values, a strictly increasing sequence drawn from them terminates; at termination
$F(\lambda_k)\le 0$, which with $\lambda_k\le\varepsilon(\chi)$ and
Proposition~\ref{prop:F} forces $\lambda_k=\varepsilon(\chi)$.
\end{proof}

Two consequences are worth stating separately, because they govern how a computed number
may be read.

\begin{remark}[One-sidedness]
Every iterate is the ratio of an explicitly exhibited coalition, so
$\lambda_k\le\varepsilon(\chi)$ holds whether or not the iteration ran to termination. A
value returned by an interrupted run is therefore a valid \emph{lower} bound on
$\varepsilon(\chi)$ and nothing more. Consequently $\lambda_k>\bar\varepsilon$ refutes
$\chi\in\mathcal{C}_{\bar\varepsilon}(v)$ even from an interrupted run, whereas
$\lambda_k\le\bar\varepsilon$ establishes membership only if the run terminated. A
nonpositive value read off a truncated measurement proves nothing.
\end{remark}

\begin{remark}[Where global optimality is needed]
The proof uses global optimality of \eqref{eq:Flambda} only in the stopping test: the
inference $F(\lambda_k)\le 0\Rightarrow\lambda_k=\varepsilon(\chi)$ fails if the reported
$F(\lambda_k)$ is merely a local or truncated maximum, since some unexplored
$(z,x)$ may still carry a larger ratio. Monotonicity and the lower-bound property survive
regardless, because they rest only on feasibility of the iterates.
\end{remark}

\begin{remark}[Sign convention]
The formulation above follows the manuscript, in which $\vmip$ is a value to be maximised
and a coalition blocks when $\vmip(S)>\ssum_{j\in S}\chi_j$. Under the cost convention, in
which values are negative for profit and a coalition blocks when it is charged more than it
would pay alone, the numerator reads $\ssum_{j\in S}\chi_j-c(S)$ and the shift in
\eqref{eq:shift} becomes $\chi'_j=\chi_j-\lambda$. The two statements are equivalent.
\end{remark}

\end{document}
